\documentclass[../../main.tex]{subfiles}
\begin{document}

{\bf [解]}

\paragraph{(1)}

$n\to\infty$の極限を考えるので$n$が十分大きいとして良い．
$y=1/x$は単調減少だから，区間$[k,k+1]$での面積比較により
\begin{align}
\sum_{k=1}^n\dfrac{1}{k}<1+\int_1^n\dfrac{1}{x}\,dx \quad\cdots\text{①}\label{eq:1} \\
\int_1^n\dfrac{1}{x}\,dx+\dfrac{1}{n+1}<\sum_{k=1}^n\dfrac{1}{k} \quad\cdots\text{②} \label{eq:2}
\end{align}
を得る．積分は
\begin{align}
 \int_1^n(1/x)dx=\log n
\end{align}
だから，\cref{eq:1,eq:2}をまとめて
\begin{align}
\log n+\dfrac{1}{n+1}<\sum_{k=1}^n\dfrac{1}{k}<\log n+1
\end{align}
を得る．両辺$\log n$（$>0$）で割って
\begin{align}
1+\dfrac{1}{(n+1)\log n}<\dfrac{1}{\log n}\sum_{k=1}^n\dfrac{1}{k}<1+\dfrac{1}{\log n} \quad\cdots\text{③} 
\end{align}
を得る．
この両辺は$n\to\infty$で$1$に収束するから，はさみうちの定理から
\begin{align}
\lim_{n\to\infty}\dfrac{1}{\log n}\sum_{k=1}^n\dfrac{1}{k}=1
\end{align}
である．

\paragraph{(2)}

与えられた関数を
\begin{align}
 f_n(x)=\displaystyle\prod_{k=0}^n(x-k) \label{eq:3}
\end{align}
とおくと積の微分法から
\begin{align}
f_n'(x)=(x-1)\cdots(x-n)+x(x-2)\cdots(x-n)+\cdots+x(x-1)\cdots(x-(n-1)) \label{eq:4}
\end{align}
となる．ここで各項は$k$番目の因子$(x-k)$を除いた積である．
$y=f_n(x)$の極値を与える$x_n$は$f_n'(x)=0$の解である．すなわち
\begin{align}
 f_n'(x_n)=(x_n-1)\cdots(x_n-n)+x_n(x_n-2)\cdots(x_n-n)+\cdots+x_n(x_n-1)\cdots(x_n-(n-1)) = 0 \label{eq:5}
\end{align}
である．
さて，\cref{eq:4}より$0\le k\le n$なる整数$k$に対し$f'_n(k)\ne 0$だから，$x_n$はこれらに該当せず，従って\cref{eq:3}より$f_n(x_n)\neq 0$だから\cref{eq:5}の両辺を$f_n(x_n)$で割って
\begin{align}
  &0=\dfrac{1}{x_n}+\dfrac{1}{x_n-1}+\cdots+\dfrac{1}{x_n-n} \\
\therefore
  &\dfrac{1}{x_n}=\dfrac{1}{1-x_n}+\dfrac{1}{2-x_n}+\cdots+\dfrac{1}{n-x_n} \label{eq:6}
\end{align}
である．これで題意の前半は示された．

次に$0<x_n\le\dfrac{1}{2}$を示す．
まず$x_n<0$と仮定すると\cref{eq:6}で左辺が負，右辺が正となって不適である．また前半の議論より$x_n \neq 0$だから$0 < x_n$である．\cref{eq:4}より
\begin{align}
 f'_n(0) &= (-1)(-2)\cdots (-n) = (-1)^{n}n! \label{eq:7}\\
 f'_n\left(\dfrac{1}{2}\right) 
   &= \left(\dfrac{1}{1/2}+\dfrac{1}{1/2-1}+\cdots+\dfrac{1}{1/2-n}\right)f_n\left(\dfrac{1}{2}\right) \\
   &= \left(\dfrac{1}{1/2}+\dfrac{1}{1/2-1}+\cdots+\dfrac{1}{1/2-n}\right)\dfrac{(-1)^{n}}{2}\prod_{k=1}^{n}\left(k-\dfrac{1}{2}\right) \label{eq:8}
\end{align}
である．$f'_n(x)$が連続であることにも注意すると，

\subparagraph{$n=1$の時}

\cref{eq:7,eq:8}より$f'_n(0)<0, f'_n(1/2)=0$より必ず$0<x \le \dfrac{1}{2}$に$f'_n(x)=0$なる$x$が存在するから
\begin{align}
 0< x_n \le \dfrac{1}{2}
\end{align}
である．

\subparagraph{$n \ge 2$の時}

この時は\cref{eq:7,eq:8}より
\begin{align}
 f'_n(0)f'_n\left(\dfrac{1}{2}\right) 
   &= (-1)^{n}n!\left(\dfrac{1}{1/2}+\dfrac{1}{1/2-1}+\cdots+\dfrac{1}{1/2-n}\right)\dfrac{(-1)^{n}}{2}\prod_{k=1}^{n}\left(k-\dfrac{1}{2}\right) \\
   &= \left(\dfrac{1}{1/2}+\dfrac{1}{1/2-1}+\cdots+\dfrac{1}{1/2-n}\right)\dfrac{n!}{2}\prod_{k=1}^{n}\left(k-\dfrac{1}{2}\right) 
\end{align}
において第一因子が負でそれ以外が正だから
\begin{align}
 f'_n(0)f'_n < 0
\end{align}
である．従って$0<x \le \dfrac{1}{2}$に$f'_n(x)=0$なる$x$が存在するから
\begin{align}
 0< x_n \le \dfrac{1}{2}
\end{align}
である．
\casesend

以上からいずれの$n$に対しても$0<x_n \le \dfrac{1}{2}$である．

\subparagraph{(3)}
$n\to\infty$を考えるので$n$が十分大きいとする．
\cref{eq:6}において$0<x_n\le \dfrac{1}{2}$より各$k\ge1$で
\begin{align}
 k - 1 < k-\dfrac{1}{2}\le k-x_n <  k
\end{align}
だから
\begin{align}
& \sum_{k=1}^n\dfrac{1}{k}<\dfrac{1}{x_n} \le \sum_{k=1}^n\dfrac{1}{k-1/2} < \dfrac{1}{1-1/2} + \sum_{k=2}^n\dfrac{1}{k-1} \\
\therefore
& \sum_{k=1}^n\dfrac{1}{k}<\dfrac{1}{x_n} < 2 + \sum_{k=1}^n\dfrac{1}{k} - \dfrac{1}{n}
\end{align}
となる．ただし2行目では右辺の和に$k-1\to k$の変換を行った．
両辺を$\log n$（$>0$）で割って
\begin{align}
\dfrac{1}{\log n}\sum_{k=1}^n\dfrac{1}{k}<\dfrac{1}{x_n\log n}<\dfrac{1}{\log n}\sum_{k=1}^n\dfrac{1}{k}+\dfrac{1}{\log n}\left(2-\dfrac{1}{n}\right) \quad\cdots\text{⑦} 
\end{align}
を得る．この両辺は(1)から$n\to\infty$で$1$に収束するので，はさみうちの定理から
\begin{align}
\lim_{n\to\infty} \dfrac{1}{x_n\log n} = 1 
\end{align}
である．$1/x$は$x>0$で連続だから
\begin{align}
\lim_{n\to\infty} x_n\log n = 1
\end{align}
が求める極限値である．

{\bf [別解]}
(2)の後半部分，$0<x_n\le\dfrac{1}{2}$は平均値の定理を用いると以下のように少し簡単に評価できる．まず解答と同じように$0<x_n$の前提で，
\begin{align}
 f_n(0) = f_n(1) = 0
\end{align}
より，平均値の定理から$0<c<1$かつ
\begin{align}
 f'_n(c) = \dfrac{f_n(1)-f_n(0)}{1-0} = 0
\end{align}
を満たす実数$c$がある．よって$0<x_n<1$である．

次に$\dfrac{1}{2}<x_n<1$と仮定すると
\begin{align}
 1<\dfrac{1}{x_n}<2, 2<\dfrac{1}{1-x_n}, 0<\dfrac{1}{2-x_n}, \cdots, 0<\dfrac{1}{n-x_n}
\end{align}
より\cref{eq:6}の左辺が$2$より小さく，右辺が$2$より大きくなって矛盾する．以上より
\begin{align}
 0 < x_n \le \dfrac{1}{2}
\end{align}
である．



\end{document}
